\documentclass[UTF8]{ctexart}
\usepackage{xeCJK}
\usepackage{geometry}
\geometry{left=2cm,right=2cm,top=2.5cm,bottom=2.5cm}
\setlength{\headheight}{12.7pt}

\usepackage{titlesec}
\titleformat{\section}
  {\normalfont\large\bfseries}
  {\thesection}
  {1em}
  {}
\titlespacing*{\section}
  {0pt}
  {3.5ex plus 1ex minus .2ex}
  {2.3ex plus .2ex}

\usepackage{amsmath}
\usepackage{amssymb}
\usepackage{amsthm}
\usepackage{enumitem}
\setlist[enumerate]{itemsep=0pt,topsep=0pt,parsep=0pt,leftmargin=0pt,itemindent=2.5em}
\setlist[itemize]{itemsep=0pt,topsep=0pt,parsep=0pt,leftmargin=0pt,itemindent=2.5em}
\usepackage{fancyhdr}
\pagestyle{fancy}
\fancyhf{}
\usepackage{graphicx}
\usepackage{hyperref}
\usepackage{indentfirst}
\usepackage{lmodern}


\begin{document}
\setlength{\abovedisplayskip}{1pt}
\setlength{\belowdisplayskip}{1pt}

\section{正规子群}

\hypertarget{sec:定义1.1}{}
\noindent\textbf{定义1.1 (正规子群)}

给定\href{01 半群与群#nameddest=sec:定义2.1}{群} $G$和它的子群$H$, 如果
\vspace{-3pt}
\[
    \forall g\in G:\quad gH=Hg,\\[-3pt]
\]
就称$H$是$G$的\textbf{正规子群}, 记作$H\lhd G$. 显然以下是等价的定义:
\begin{enumerate}
    \item 对任意的$g\in G$, 有$gHg^{-1}=H$.
    \item 对任意的$g\in G$和$h\in H$, 有$ghg^{-1}\in H$.
    \item 对任意的$\varphi\in\mathrm{Int}(G)$, 有$\varphi[H]=H$.
\end{enumerate}

\vspace{10pt}

\hypertarget{sec:定理1.1}{}
\noindent\textbf{定理1.1}

任给群$K<H<G$, 如果$K\lhd G$, 那么$K\lhd H$.

\vspace{10pt}

显然Abel群的子群都是正规子群. 另外, 群同态的核也是正规子群.

\vspace{10pt}

\hypertarget{sec:定理1.2}{}
\noindent\textbf{定理1.2}

任给群同态$\varphi:G\to G'$, 它的核$\mathrm{Ker}\,\varphi$是$G$的正规子群.

\noindent{\kaishu 证明.}

对任意的$g\in G$和$h\in\mathrm{Ker}\,\varphi$, 有
\[
    \varphi(h)=1_{G'}\quad\implies\quad
    \varphi(ghg^{-1})=\varphi(g)\cdot1_{G'}\cdot\varphi(g)^{-1}=1_{G'},
\]
即$ghg^{-1}\in\mathrm{Ker}\,\varphi$. \hfill $\qedsymbol$

\vspace{10pt}

\hypertarget{sec:定义1.2}{}
\noindent\textbf{定义1.2 (群的中心)}

任给群$G$, 定义其\textbf{中心}为
\[
    \mathrm{Z}(G)\overset{\mathrm{def}}{=}\{g\in G\mid\forall h\in G:gh=hg\},
\]
它恰好是$\mathrm{Int}$的核, 从而是$G$的正规子群.

\noindent{\kaishu 验证.}

对任意的$g\in G$, $\mathrm{Int}(g)=\mathrm{id}_{G}$当且仅当
\[
    \forall h\in G:ghg^{-1}=h,
\]
即$g\in \mathrm{Z}(G)$. \hfill $\qedsymbol$





\newpage

\section{商群}

\hypertarget{sec:定义2.1}{}
\noindent\textbf{定义2.1 (商群)}

给定群$G$和它的子群$H$, 定义
\[
    \,\!^{\displaystyle G}\!/_{\displaystyle H}\overset{\mathrm{def}}{=}
    \{gH\mid g\in G\},
\]
如果$H\lhd G$,
则$\left(\,\!^{\displaystyle G}\!/_{\displaystyle H},H,\cdot\right)$
是群, 称之为$G$对$H$的\textbf{商群}.

\noindent{\kaishu 验证.}

对任意的$g_1,g_2\in G$, 注意到$g_1Hg_2H=g_1g_2HH=g_1g_2H$, 从而
\vspace{3pt}
\begin{enumerate}
    \item $\,\!^{\displaystyle G}\!/_{\displaystyle H}$对乘法封闭,
    \vspace{3pt}
    \item 对任意的$gH\in\,\!^{\displaystyle G}\!/_{\displaystyle H}$,
    有$gH\cdot H=H\cdot gH=H$,
    \vspace{3pt}
    \item 对任意的$gH\in\,\!^{\displaystyle G}\!/_{\displaystyle H}$,
    有$gH\cdot g^{-1}H=g^{-1}H\cdot gH=H$. \hfill $\qedsymbol$
\end{enumerate}

\vspace{10pt}

\hypertarget{sec:定义2.2}{}
\noindent\textbf{定义2.2 (投影同态)}

任给群$G$和它的正规子群$H$, 有群同态
\[
    \pi:G\to\,\!^{\displaystyle G}\!/_{\displaystyle H},\quad g\mapsto gH,\\[3pt]
\]
称之为\textbf{投影同态}, 并且$\mathrm{Ker}\,\pi=H$,
$\mathrm{Im}\,\pi=\,\!^{\displaystyle G}\!/_{\displaystyle H}$.

\vspace{15pt}

\hypertarget{sec:定理2.1}{}
\noindent\textbf{定理2.1 (商群的泛性质)}

任给群同态$\varphi:G\to G'$, 如果有$H\lhd G$且$H<\mathrm{Ker}\,\varphi$,
那么存在唯一的群同态
\vspace{3pt}
\[
    \overline{\varphi}:\,\!^{\displaystyle G}\!/_{\displaystyle H}\to G'\\[3pt]
\]
使得$\varphi=\overline{\varphi}\circ\pi$, 并且
$\mathrm{Ker}\,\overline{\varphi}=
\,\!^{\displaystyle\mathrm{Ker}\,\varphi}\!/_{\displaystyle H}$,
$\mathrm{Im}\,\overline{\varphi}=\mathrm{Im}\,\varphi$.

\noindent{\kaishu 证明.}

对任意的$g_1,g_2\in G$, 如果$\pi(g_1)=\pi(g_2)$, 那么
\vspace{3pt}
\[
    g_1H=g_2H\implies g_2^{-1}g_1\in H\subset\mathrm{Ker}\,\varphi\implies
    \varphi(g_2^{-1}g_1)=1_{G'}\implies\varphi(g_1)=\varphi(g_2),\\[3pt]
\]
从而存在唯一的映射$\overline{\varphi}:\,\!^{\displaystyle G}\!/_{\displaystyle H}
\to G'$使得$\varphi=\overline{\varphi}\circ\pi$, 即对任意的$g\in G$有
$\overline{\varphi}(gH)=\varphi(g)$.

此时对任意的$g_1,g_2\in G$, 有
\[
    \overline{\varphi}(g_1g_2H)=\varphi(g_1g_2)=\varphi(g_1)\cdot\varphi(g_2)=
    \overline{\varphi}(g_1H)\cdot\overline{\varphi}(g_2H),\\[3pt]
\]
从而$\overline{\varphi}:\,\!^{\displaystyle G}\!/_{\displaystyle H}
\to G'$是群同态. 最后,
\vspace{5pt}
\begin{equation*}\left\{\begin{aligned}
    &\mathrm{Ker}\,\overline{\varphi}=\{gH\mid\overline{\varphi}(gH)=1_{G'}\}=
    \{gH\mid\varphi(g)=1_{G'}\}=\{gH\mid g\in\mathrm{Ker}\,\varphi\}=\,
    \,\!^{\displaystyle\mathrm{Ker}\,\varphi}\!/_{\displaystyle H},\\[3pt]
    &\mathrm{Im}\,\overline{\varphi}=\{\overline{\varphi}(gH)\mid g\in G\}=
    \{\varphi(g)\mid g\in G\}=\mathrm{Im}\,\varphi.
\end{aligned}\right.\tag*{\qedsymbol}\end{equation*}

\end{document}